% !TEX root = ../main.tex
\chapter{Stochastic Processes I: Random Walks, Markov Chains, Martingales}\label{ch:sp1}
\begin{paracol}{2}

\begin{sources}
\begin{tabularx}{\linewidth}{@{}l l L@{}}
\toprule
\textbf{Lecture} & \textbf{Speaker} & \textbf{Content}\\
\midrule
2024 L5 (from 1:12:25) & P.~Kempthorne & Definition of a martingale; the simple random walk and simulations of its paths; the martingale $S_n^2-n\sigma^2$.\\
2024 L6 (first 67 min) & P.~Kempthorne & Product and exponential martingales; filtrations, non-anticipating strategies and martingale transforms; stopping times; gambler's ruin (fair and biased) by optional stopping; Markov processes and chains, $n$-step transitions; credit-rating migration; an up/down Markov chain for AAPL.\\
2013 L5 & C.~Lee & Processes as distributions over paths; simple random walk, its properties, gambler's ruin by recursion; Markov chains, transition matrix, stationary distribution via Perron--Frobenius; martingales, stopping times, optional stopping theorem.\\
\bottomrule
\end{tabularx}

\smallskip
\textbf{Files.} Slides \file{mit18_642_f24_lec05.pdf} (\emph{Stochastic Processes I}, 44 slides in 88 pages of
incremental builds); 2013 lecture notes \file{MIT18_S096F13_lecnote5.pdf}. The R script
\file{MC_Example1.r} used for the AAPL example (slide 44) was distributed on Canvas only and is not among
the OCW files.
\textbf{Problem sets.} 2024 PS2 Problem~1, PS3 Problem~3; 2013 PS2 Problems A-3, A-4, B-5.
\textbf{Elsewhere.} The first 72 minutes of 2024 L5 (PCA, $\chi^2$/$t$/$F$ distributions, LLN, CLT,
a one-period CAPM) are in \cref{ch:probability}; the second half of 2024 L6 (regression) is in
\cref{ch:regression}.
\end{sources}

\section*{Motivation}
A price history is one realisation of a \emph{stochastic process}. Standing at time $t$ we know the
past path and want to say something about the future: whether it depends on the past (can prices be
predicted?), how the process behaves in the long run, and how likely extreme ``boundary'' events are
\ts{13}{5}{8:25}. This chapter introduces three families of discrete-time processes that answer these
questions in increasing generality and that return, in continuous time, for the rest of the course:
\begin{itemize}
  \item the \emph{random walk}, the simplest model of a price and the discrete skeleton of Brownian
        motion (\cref{ch:sp2});
  \item \emph{Markov chains}, where the future depends on the past only through the present state
        (credit ratings, market regimes); they extend the stochastic matrices of \cref{sec:markov};
  \item \emph{martingales}, the mathematical model of a fair game. Discounted prices are martingales
        under the risk-neutral measure, gains from trading are martingale transforms (discrete
        stochastic integrals, \cref{ch:stochcalc}), and the optional stopping theorem turns
        hard probability questions into one-line computations.
\end{itemize}
For a physicist: the random walk is diffusion, a Markov chain is a master equation relaxing to
equilibrium, and a martingale is a quantity that is conserved \emph{on average}; optional stopping is the
use of such a conservation law to solve a boundary-value problem.

% ------------------------------------------------------------------
\section{Stochastic processes}\label{ch04:sec-processes}

\begin{definition}[Stochastic process]
A \emph{stochastic process} is a family $\{X_t,\ t\in\mathcal T\}$ of random variables on a common
probability space $(\Omega,\mathcal F,\Prob)$, indexed by \emph{time}: $\mathcal T=\{0,1,2,\dots\}$
(discrete time) or $\mathcal T=[0,\infty)$ (continuous time) \ts{13}{5}{1:03}. Equivalently, each outcome
$\omega\in\Omega$ selects a whole \emph{path} $t\mapsto X_t(\omega)$, so a process is a probability
distribution on a space of paths \ts{13}{5}{3:09}.
\end{definition}
The first reading is the intuitive one (a random value at each time); the second is the one needed to do
mathematics: at time $t$ we observe the initial segment of a path drawn once and for all, and we reason
about the unobserved rest of it.

\begin{example}[Same marginals, different processes {\ts{13}{5}{4:10}}]\label{ch04:ex-threeproc}
For $t\ge0$ consider
\begin{enumerate}[(a)]
  \item $X_t=t$ with probability~1;
  \item $X_t=\varepsilon t$ for all $t$, with a
      \emph{single} sign $\varepsilon=\pm1$ drawn with probability $1/2$ each;
  \item $X_t=\varepsilon_tt$ with an
      independent sign $\varepsilon_t$ for every $t$.
\end{enumerate}
In (b) and (c) every $X_t$ has the same distribution
($\pm t$ with probability $1/2$), but in (b) one observation reveals the entire future, while in (c) the
past says nothing about the future beyond $X_t\in\{t,-t\}$ \ts{13}{5}{7:20}. A process is \emph{not}
determined by its one-time marginals: what matters is the joint law of the path.
\end{example}

% ------------------------------------------------------------------
\section{The random walk}\label{ch04:sec-rw}

\begin{definition}[Random walk]\label{ch04:def-rw}
Let $X_1,X_2,\dots$ be i.i.d.\ with mean $\mu$ and variance $\sigma^2$, and let $S_0=0$,
$S_n=X_1+\dots+X_n$. $\{S_n\}$ is a \emph{random walk} with steps $X_i$. When
$\Prob(X_i=+1)=\Prob(X_i=-1)=\tfrac12$ it is the (one-dimensional) \emph{simple random walk} (SRW)
\ts{13}{5}{12:35}, \ts{24}{5}{1:14:36}.
\end{definition}

The path can be drawn as the graph of $n\mapsto S_n$, or as the trajectory of a walker taking unit steps
left or right on a line \ts{13}{5}{13:40}. The 2013 course and the 2024 problem sets also use the
\emph{biased} walk ($\Prob(X=+1)=p\ne\tfrac12$) and walks with a zero step (2024 PS2 Problem~1,
\cref{ch04:ex-ps2-1}).

\begin{proposition}[Basic properties]\label{ch04:prop-rw}
For a random walk:
\begin{enumerate}[(i)]
  \item $\E[S_n]=n\mu$ and $\Var(S_n)=n\sigma^2$; for the SRW $\E[S_n]=0$, $\Var(S_n)=n$ \ts{24}{5}{1:15:36};
  \item (independent increments) for $0\le k_0\le k_1\le\dots\le k_r$ the increments
        $S_{k_1}-S_{k_0},\dots,S_{k_r}-S_{k_{r-1}}$ are mutually independent \ts{13}{5}{20:56};
  \item (stationary increments) $S_{k+h}-S_k\eqd S_h$ for all $k,h\ge0$ \ts{13}{5}{22:03};
  \item (CLT) $(S_n-n\mu)/(\sigma\sqrt n)\to N(0,1)$ in distribution.
\end{enumerate}
\end{proposition}
\begin{proof}
(i) Linearity of expectation, and additivity of variance for independent summands.

(ii) The increment over $(k_{i-1},k_i]$ is a function of the block $X_{k_{i-1}+1},\dots,X_{k_i}$;
functions of disjoint blocks of independent variables are independent.

(iii) $S_{k+h}-S_k=X_{k+1}+\dots+X_{k+h}$ is a sum of $h$ i.i.d.\ copies of $X$, like $S_h$.

(iv) Central limit theorem, \cref{ch:probability}.
\end{proof}
Properties (ii) and (iii) are exactly those that define Brownian motion in continuous time
\ts{13}{5}{22:03}: the random walk is ``the fundamental stochastic process'', and much of
\cref{ch:sp2} is its continuous limit.

\paragraph{The $\sqrt n$ scale.} By (iv) the SRW at time $n$ is typically at distance of order
$\sqrt n$ from the origin, not $n$: although $|S_n|$ could be as large as $n$, the probability that
$|S_n|\le 2\sqrt n$ tends to $0.954$ \ts{13}{5}{17:51}. Simulated paths (slides 4--8) become more
jagged as $n$ grows and stay mostly inside the envelopes $\pm\sqrt n$, $\pm2\sqrt n$, $\pm3\sqrt n$
\ts{24}{5}{1:16:43}; see \cref{ch04:fig-rw}. Nevertheless the envelopes are crossed again and again: with
probability one the SRW returns to $0$ infinitely often, and for every $c>0$ it exceeds $c\sqrt n$
infinitely often \ts{13}{5}{18:52} (a consequence of the law of the iterated logarithm,
$\limsup_n S_n/\sqrt{2n\ln\ln n}=1$ a.s., not proved in the course).

\begin{figure}[ht]
\centering
\begin{tikzpicture}[x=0.058cm,y=0.075cm,>=Stealth]
  \draw[black!40] (0,-40) rectangle (200,40);
  \draw[black!25] (0,0) -- (200,0);
  \foreach \yy in {-40,-20,0,20,40} \draw[black!40] (0,\yy) -- (-2,\yy) node[left,font=\scriptsize,black]{$\yy$};
  \foreach \xx in {0,50,100,150,200} \draw[black!40] (\xx,-40) -- (\xx,-42) node[below,font=\scriptsize,black]{$\xx$};
  \node[below,font=\small] at (100,-47) {step $n$};
  \node[rotate=90,font=\small] at (-17,0) {$S_n$};
  \draw[thmcol,thick,domain=0:200,samples=100] plot (\x,{sqrt(\x)});
  \draw[thmcol,thick,domain=0:200,samples=100] plot (\x,{-sqrt(\x)});
  \draw[thmcol,thick,dashed,domain=0:200,samples=100] plot (\x,{2*sqrt(\x)});
  \draw[thmcol,thick,dashed,domain=0:200,samples=100] plot (\x,{-2*sqrt(\x)});
  \pgfmathsetseed{3}
  \foreach \n [evaluate=\n as \s using {\p+2*round(rnd)-1}, remember=\s as \p (initially 0)] in {1,...,200}
    {\draw[defcol] (\n-1,\p) -- (\n,\s);}
  \pgfmathsetseed{17}
  \foreach \n [evaluate=\n as \s using {\p+2*round(rnd)-1}, remember=\s as \p (initially 0)] in {1,...,200}
    {\draw[intcol] (\n-1,\p) -- (\n,\s);}
  \pgfmathsetseed{29}
  \foreach \n [evaluate=\n as \s using {\p+2*round(rnd)-1}, remember=\s as \p (initially 0)] in {1,...,200}
    {\draw[fincol] (\n-1,\p) -- (\n,\s);}
  \pgfmathsetseed{41}
  \foreach \n [evaluate=\n as \s using {\p+2*round(rnd)-1}, remember=\s as \p (initially 0)] in {1,...,200}
    {\draw[cscol] (\n-1,\p) -- (\n,\s);}
  \node[thmcol,font=\scriptsize,anchor=west] at (201,{sqrt(200)}) {$\sqrt n$};
  \node[thmcol,font=\scriptsize,anchor=west] at (201,{2*sqrt(200)}) {$2\sqrt n$};
  \node[thmcol,font=\scriptsize,anchor=west] at (201,{-sqrt(200)}) {$-\sqrt n$};
  \node[thmcol,font=\scriptsize,anchor=west] at (201,{-2*sqrt(200)}) {$-2\sqrt n$};
\end{tikzpicture}
\caption{Four simulated paths of the simple random walk (200 steps) with the one- and two-standard-deviation
envelopes $\pm\sqrt n$, $\pm2\sqrt n$; slides 4--8 of \file{lec05} show the same picture for up to 10\,000
steps \ts{24}{5}{1:16:43}.}\label{ch04:fig-rw}
\end{figure}

\begin{coursenote}[Slips in the 2013 lecture]
At \ts{13}{5}{16:47} the value $S_t$ is said to be approximately normal ``with mean 0 and variance
$\sqrt t$''; the variance is $t$ and the \emph{standard deviation} is $\sqrt t$, which is the scale meant
in the rest of the discussion. (The $1/t$ proposed by a student is the variance of the average $S_t/t$.)
\end{coursenote}

\begin{intuition}[Diffusion]
$S_n$ is the position of a particle receiving independent unit kicks: $\langle S_n^2\rangle=n$ is
Einstein's $\langle x^2\rangle=2Dt$ with $D=\tfrac12$ (unit step, unit time). Rescaling space by
$\sqrt n$ and time by $n$, $t\mapsto S_{\lfloor nt\rfloor}/\sqrt n$ converges to Brownian motion
(Donsker's theorem, \cref{ch:sp2}); the jagged look of the long simulated paths is the nowhere-differentiability
of Brownian paths in the making.
\end{intuition}

\begin{finance}[Random walks as price models]
L.~Bachelier (1900), cited in the 2013 notes, modelled prices on the Paris exchange as a random walk;
the 2013 notes call the SRW ``a (rather inaccurate) model of stock price''. The course refines it in
two directions that appear in the problem sets: a drift, calibrated to an expected daily price change
(2024 PS2 Problem~1, where an intraday price moves in 78 five-minute steps), and a random walk for
the \emph{logarithm} of the price, $\log P_t=\log P_{t-1}+X_t$ (2024 PS3 Problem~3), which keeps prices
positive and makes returns, not price changes, i.i.d. Time-varying volatility comes later
(\cref{ch:volatility}).
\end{finance}

\begin{exercise}[PS2 (2024), Problem 1: random walk with bias]\label{ch04:ex-ps2-1}
Let $X_n$ be i.i.d.\ steps with $\Prob(X_n=+1)=p$, $\Prob(X_n=0)=1-(p+q)$, $\Prob(X_n=-1)=q$, where $p,q>0$ and $p+q<1$
(symmetric if $p=q$, biased upwards if $p>q$), and $S_n=\sum_{i=1}^nX_i$, $S_0=0$.
\begin{enumerate}[(a)]
  \item Derive the moment generating function $M_{X_i}(t)$ of a step and the mean and variance of $X_i$.
  \item Derive the moment generating function $M_{S_n}(t)$ and the mean and variance of $S_n$.
  \item With $\mu$, $\sigma^2$ the mean and variance of a step, prove that $Z_n=(S_n-n\mu)/\sqrt{n\sigma^2}$ converges in
      distribution to $N(0,1)$.
  \item Model the intraday price of a stock as such a walk with 78 five-minute steps (the six and a half hours of NYSE
      trading), with $p=0.52$ and $q=0.47$. Which step size gives an expected daily price change of $+\$0.50$? With that
      step size, what is the standard deviation of the daily price change?
\end{enumerate}
\end{exercise}

\begin{exercise}[\file{lecnote5}, Proposition 2.1]
Prove the independent- and stationary-increment properties of the random walk (\cref{ch04:prop-rw}) directly from
the definition, and check which of them survive if the steps are independent but not identically distributed.
\end{exercise}

% ------------------------------------------------------------------
\section{Markov chains}\label{ch04:sec-mc}

\subsection{The Markov property}
\begin{definition}[Markov process, Markov chain]\label{ch04:def-markov}
A process $\{X_t\}$ is a \emph{Markov process} if for all times $u<s<t$, given the present value $X_s$,
the future value $X_t$ is independent of the past values $X_u$ \ts{24}{6}{36:30}:
\[
  \bigl[X_t\mid X_r,\ r\le s\bigr]\ \eqd\ \bigl[X_t\mid X_s\bigr].
\]
A discrete-time Markov process with a finite or countable \emph{state space} $S$ is a \emph{Markov chain}:
\begin{equation}\label{ch04:eq-markov}
  \Prob(X_{n+1}=j\mid X_0=i_0,\dots,X_{n-1}=i_{n-1},X_n=i)=\Prob(X_{n+1}=j\mid X_n=i)
\end{equation}
for all states and all $n$ \ts{24}{6}{39:38}.
\end{definition}
``The effect of the past on the future is summarised by the current state'' \ts{13}{5}{33:32}: two paths
arriving at the same state at time $s$ have the same future distribution, however they got there
\ts{24}{6}{38:30}.

\begin{example}[Random walks are Markov]
If $S_{n+1}=S_n+X_{n+1}$ with $X_{n+1}$ independent of $(S_0,\dots,S_n)$, then
$\Prob(S_{n+1}\in\cdot\mid S_0,\dots,S_n)$ is the law of $S_n+X$, a function of $S_n$ only. For the SRW,
$\Prob(S_{n+1}=s\mid S_0,\dots,S_n)=\tfrac12$ if $s=S_n\pm1$ and $0$ otherwise \ts{13}{5}{37:45}. The
state space $\mathbb Z$ is infinite, so there is no finite transition matrix \ts{13}{5}{43:12}.
\end{example}

\begin{exercise}[PS2 (2013), A-3: identify Markov processes]
Which of the following are Markov processes?
\begin{enumerate}[(a)]
  \item The simple random walk $S_t$.
  \item $X_t=|S_t|$.
  \item $X_0=0$,
      $X_{t+1}=X_t+(Y_t-\tfrac1\lambda)$ with $Y_t$ i.i.d.\ exponential with parameter $\lambda$.
  \item $X_0=0$,
      $X_{t+1}=X_t+Z_tZ_{t-1}\cdots Z_0$ with $Z_t$ i.i.d.\ log-normal.
  \item $X_0=0$, $X_{t+1}=X_t+W_t+W_{t-1}+\dots+W_0$ with
      $W_t$ i.i.d.\ normal.
\end{enumerate}
\end{exercise}

\subsection{Transition matrices}
The one-step transition probabilities $\Prob(X_{n+1}=j\mid X_n=i)$ together with the initial
distribution $p_i=\Prob(X_0=i)$ specify the law of the whole process \ts{24}{6}{41:45}. The chain has
\emph{stationary} (time-homogeneous) transition probabilities if these do not depend on $n$
\ts{24}{6}{42:52}; then they form the \emph{transition matrix}
\[
  P=(P_{ij})_{i,j\in S},\qquad P_{ij}=\Prob(X_{n+1}=j\mid X_n=i),\qquad P_{ij}\ge0,\quad \sum_{j}P_{ij}=1,
\]
a \emph{row}-stochastic matrix (row $i$ = where you are, column $j$ = where you go) \ts{24}{6}{43:58}.
All chains below are time-homogeneous.

\begin{coursenote}[Row versus column conventions]
The 2024 slides use the row-stochastic $P$ above, acting on
\emph{row} vectors of probabilities from the right. The 2013 notes (and
\cref{sec:markov}) use the column-stochastic $A=P\T$ acting on column
vectors: $\bm\pi(t+1)=A\bm\pi(t)$ is the transpose of
$\bm p^{(t+1)}=\bm p^{(t)}P$. This chapter uses $P$. In the 2013
lecture the machine example below was first written in the row form
and then ``flipped'' \ts{13}{5}{51:41}.
\end{coursenote}

\begin{proposition}[Law of the path]\label{ch04:prop-path}
$\Prob(X_0=i_0,X_1=i_1,\dots,X_n=i_n)=p_{i_0}\,P_{i_0i_1}P_{i_1i_2}\cdots P_{i_{n-1}i_n}$.
\end{proposition}
\begin{proof}
Write $A=\{X_0=i_0,\dots,X_{n-1}=i_{n-1}\}$ and $B=\{X_n=i_n\}$. Then $\Prob(A\cap B)=\Prob(A)\Prob(B\mid A)$
(joint = marginal $\times$ conditional), $\Prob(B\mid A)=\Prob(X_n=i_n\mid X_{n-1}=i_{n-1})$ by the Markov
property, and this equals $P_{i_{n-1}i_n}$ by stationarity. Induction on $n$ \ts{24}{6}{45:04}, \ts{24}{6}{48:09}.
\end{proof}

\begin{theorem}[$n$-step transitions, Chapman--Kolmogorov]\label{ch04:thm-nstep}
Let $P^{(n)}_{ij}=\Prob(X_{m+n}=j\mid X_m=i)$ (independent of $m$). Then
\[
  P^{(n)}_{ij}=\sum_{k\in S}P_{ik}\,P^{(n-1)}_{kj},\quad P^{(0)}=I,\qquad\text{hence}\qquad P^{(n)}=P^n,
  \quad P^{(m+n)}=P^{(m)}P^{(n)} .
\]
The distribution of $X_n$ is the row vector $\bm p^{(n)}=\bm p\,P^n$, i.e.\
$\Prob(X_n=k)=\sum_jp_jP^{(n)}_{jk}$ \ts{24}{6}{52:29}.
\end{theorem}
\begin{proof}[Proof (first-step analysis)]
Condition on the first step \ts{24}{6}{49:16}:
\begin{align*}
P^{(n)}_{ij}&=\sum_k\Prob(X_1=k\mid X_0=i)\,\Prob(X_n=j\mid X_1=k,X_0=i)
 =\sum_kP_{ik}\,\Prob(X_n=j\mid X_1=k)\\
 &=\sum_kP_{ik}\,\Prob(X_{n-1}=j\mid X_0=k)=\sum_kP_{ik}P^{(n-1)}_{kj},
\end{align*}
using the Markov property and then stationarity. In matrix form
$P^{(n)}=PP^{(n-1)}$ \ts{24}{6}{51:21}; induction gives $P^n$.
\end{proof}
The 2013 notes decompose on the \emph{last} step instead, $P^{(n)}_{ij}=\sum_kP^{(n-1)}_{ik}P_{kj}$, and
the lecture reads the two-step matrix as a matrix product (``sum over all intermediate states'')
\ts{13}{5}{42:08}. First-step analysis is the more useful habit: it is how hitting probabilities and
expected hitting times are computed (\cref{ch04:sec-ruin}).

\begin{example}[A machine that breaks {\ts{13}{5}{44:15}}]\label{ch04:ex-machine}
A machine is working (W) or broken (B) each day. If working it breaks the next day with probability
$0.01$; if broken it is repaired by the next day with probability $0.8$:
\[
  P=\begin{pmatrix}0.99&0.01\\0.8&0.2\end{pmatrix},\qquad
  P^2=\begin{pmatrix}0.9881&0.0119\\0.952&0.048\end{pmatrix}.
\]
What is the state distribution after ten years, $\bm e_W P^{3650}$ \ts{13}{5}{46:26}? The lecture
``cheats'': after so long the distributions on day 3650 and 3651 should coincide, so the answer $\bm\pi$ should
satisfy $\bm\pi=\bm\pi P$ \ts{13}{5}{47:27}. This is the two-state chain of \cref{ex:twostate} with
$p=0.01$, $q=0.8$: $\bm\pi=(0.8,\,0.01)/0.81=(0.98765,\,0.01235)$ and the second eigenvalue is
$1-p-q=0.19$, so $P^n$ converges to the matrix with both rows equal to $\bm\pi$ like $0.19^n$ and
$P^{3650}$ equals it to machine precision. In the long run the machine is broken on $1.23\%$ of the days.
\end{example}

\subsection{Stationary distributions and convergence}
\begin{definition}[Stationary distribution]
A probability (row) vector $\bm\pi$ on $S$ is \emph{stationary} if $\bm\pi P=\bm\pi$, i.e.\
$\pi_j=\sum_k\pi_kP_{kj}$ for all $j$ \ts{13}{5}{53:49}. If $X_0\sim\bm\pi$ then $X_n\sim\bm\pi$ for every
$n$ \ts{13}{5}{54:52}. Equivalently, $\bm\pi\T$ is an eigenvector of eigenvalue $1$ of the column-stochastic
$A=P\T$.
\end{definition}

\begin{theorem}[Convergence to equilibrium]\label{ch04:thm-conv}
Let $S$ be finite with $m$ states and suppose some power $P^k$ has all entries positive (in particular, if
all $P_{ij}>0$). Then there is a unique stationary distribution $\bm\pi$, all $\pi_j>0$, and for every
initial distribution
\[
  \bigl\|\bm p\,P^n-\bm\pi\bigr\|_1\le 2\,(1-m\delta)^{\lfloor n/k\rfloor},\qquad
  \delta=\min_{i,j}(P^k)_{ij}>0 ;
\]
in particular $P^{(n)}_{ij}\to\pi_j$ for all $i,j$: the chain forgets its initial state geometrically fast.
\end{theorem}

\paragraph{The course's argument} (2013 notes, Theorem~3.3, for $k=1$ \ts{13}{5}{49:36}).
$A=P\T$ is column-stochastic with positive entries. By \cref{prop:stochastic}
its eigenvalues satisfy $|\lambda|\le1$ and $1$ is one of them; by
Perron--Frobenius (\cref{thm:pf}) the dominant eigenvalue is simple, has a
positive eigenvector, and all other eigenvalues have $|\lambda|<1$. The positive
eigenvector, normalised, is $\bm\pi$, and the other modes of $A^n\bm p\T$
decay (as in the state equations of \cref{sec:eigen}), so
$A^n\bm p\T\to\bm\pi\T$ \ts{13}{5}{52:41}.

\begin{proof}[A self-contained proof (not from the course)]
For two probability vectors $\bm\mu,\bm\nu$ let $\bm d=\bm\mu-\bm\nu$, so
$\sum_id_i=0$. Then, writing $Q=P^k$,
\begin{align*}
  (\bm dQ)_j&=\sum_id_iQ_{ij}=\sum_id_i\,(Q_{ij}-\delta)\\
  &\quad\Rightarrow\quad
  \|\bm dQ\|_1\le\sum_i|d_i|\sum_j(Q_{ij}-\delta)=(1-m\delta)\,\|\bm d\|_1,
\end{align*}
because $Q_{ij}-\delta\ge0$ and each row of $Q$ sums to $1$; similarly
$\|\bm dP\|_1\le\|\bm d\|_1$. Hence $\bm\mu\mapsto\bm\mu Q$ is a contraction
of the (complete) simplex in the $\ell^1$ distance, and by Banach's
fixed-point theorem has a unique fixed point $\bm\pi$. Since
$(\bm\pi P)Q=(\bm\pi Q)P=\bm\pi P$, $\bm\pi P$ is also a fixed point of
$Q$, so $\bm\pi P=\bm\pi$; any stationary distribution of $P$ is fixed
by $Q$, hence equals $\bm\pi$. Finally
\begin{align*}
  \|\bm pP^n-\bm\pi\|_1&=\|(\bm p-\bm\pi)P^n\|_1\\
  &\le(1-m\delta)^{\lfloor n/k\rfloor}\|\bm p-\bm\pi\|_1,
\end{align*}
with $\|\bm p-\bm\pi\|_1\le2$, and $\pi_j=(\bm\pi Q)_j\ge\delta>0$.
\end{proof}

\begin{intuition}[Relaxation to equilibrium]
$\bm p^{(n+1)}=\bm p^{(n)}P$ is a discrete-time master equation. The stationary distribution is its
equilibrium and $1-|\lambda_2|$ (the spectral gap) is the inverse relaxation time, as in \cref{ex:twostate}.
The contraction proof has a probabilistic meaning (coupling): run two copies of the chain from different
starts; at every block of $k$ steps they have probability at least $m\delta$ of landing in the same state,
after which they can be glued together. The initial condition is forgotten as soon as the copies have
met.
\end{intuition}

\begin{coursenote}[The puzzle left open in the 2013 lecture]
For the machine example the lecture sums the two components of $A\bm v=\lambda\bm v$ ($A$ column-stochastic)
and gets $\sum_iv_i=\lambda\sum_iv_i$, hence $\lambda=1$; then notices that the same argument seems to
force \emph{every} eigenvalue to be $1$, which is absurd, and leaves the point open \ts{13}{5}{56:57}. The
resolution: the identity only says that $\lambda=1$ \emph{or} $\sum_iv_i=0$. Eigenvectors of the other
eigenvalues have components summing to zero (like $(1,-1)\T$ in \cref{ex:twostate}), so they cannot be
normalised to probability vectors. For the Perron--Frobenius eigenvector, which is positive, the sum is
non-zero and the argument correctly gives $\lambda_0=1$.
\end{coursenote}

\paragraph{When the hypotheses fail.}
\begin{itemize}
  \item \emph{Periodicity.} The 2013 notes (Example~3.2) consider the walk $X_{n+1}=X_n\pm1\pmod N$ on the
        cycle $\mathbb Z_N$, with stationary distribution $\pi_i=1/N$ (the matrix is doubly stochastic, 2013 PS2
        A-4). If $N$ is even the chain alternates between even and odd states, no power of $P$ is positive and
        $P^{(n)}_{ij}$ does not converge: a stationary distribution need not be a limit (compare the chain
        $\bigl(\begin{smallmatrix}0&1\\1&0\end{smallmatrix}\bigr)$ in \cref{sec:markov}).
  \item \emph{Reducibility.} If the state space splits into pieces that do not communicate, the stationary
        distribution need not be unique (2013 PS2 A-4(a)). States $i,j$ \emph{communicate} if each can be
        reached from the other in some number of steps; a chain whose transition graph is strongly connected
        is \emph{irreducible} \ts{24}{6}{1:03:03}. For a finite irreducible chain the stationary distribution is
        unique and positive; if moreover it is aperiodic, some $P^k$ is positive and \cref{ch04:thm-conv}
        applies (standard facts, not proved in the course).
  \item \emph{Infinite state space.} ``A corresponding theorem is not true'' (2013 notes): the SRW on $\mathbb Z$
        has no stationary distribution. Indeed $\pi_j=\tfrac12(\pi_{j-1}+\pi_{j+1})$ makes $j\mapsto\pi_j$ linear;
        being non-negative on all of $\mathbb Z$ it is constant, and a constant cannot sum to $1$.
\end{itemize}

\paragraph{Estimating a chain from data (not discussed in class).} Given an observed path with $n_{ij}$
transitions $i\to j$, \cref{ch04:prop-path} gives the likelihood $\prod_{i,j}P_{ij}^{\,n_{ij}}$ (times
$p_{i_0}$). Maximising $\sum_{ij}n_{ij}\log P_{ij}$ subject to $\sum_jP_{ij}=1$ with Lagrange multipliers gives
the empirical frequencies $\hat P_{ij}=n_{ij}/\sum_kn_{ik}$; this is what the R package \texttt{markovchain}
uses in the case study below.

\begin{exercise}[PS2 (2013), A-4]
\leavevmode
\begin{enumerate}[(a)]
  \item Construct a finite Markov chain whose stationary distribution is not unique.
  \item A chain is \emph{doubly stochastic} if, in addition to $\sum_jp_{ij}=1$, also $\sum_ip_{ij}=1$ for every $j$ (the
      problem phrases this as ``each row sums to one'' for the transposed matrix $A$ of the 2013 notes). Prove that the uniform
      distribution $\pi_j=1/m$ is stationary for a doubly stochastic chain on $m$ states.
\end{enumerate}
\end{exercise}

\begin{exercise}[PS2 (2013), B-5: walk with reflecting barriers]
A person walks on the line, stepping right with probability $b$ and left with probability $1-b$ at each period. He
starts in one of the positions $1,\dots,m$; a step to $0$ (or to $m+1$) is instantly reflected back to $1$ (or to $m$).
Equivalently, in position $1$ (position $m$) he stays put with probability $1-b$ (probability $b$).
\begin{enumerate}[(a)]
  \item Model this as a Markov chain and write its transition matrix.
  \item Find the stationary distribution.
\end{enumerate}
\end{exercise}

\subsection{Markov chains in finance}

\paragraph{Credit ratings.} Rating agencies (S\&P, Moody's, Fitch) grade corporate and sovereign borrowers
from AAA (highest) through AA, A, BBB (investment grade) to BB, B, CCC, CC, C (speculative grade) and D
(default) \ts{24}{6}{52:29}. The rating determines the interest rate at which a company can borrow by
issuing bonds \ts{24}{6}{53:29}. Modelling the rating of an issuer as a Markov chain with a one-year
\emph{migration matrix}, default being an absorbing state ($P_{DD}=1$), the probability of default within
$n$ years for an issuer rated $i$ today is $(P^n)_{iD}$. Slide~43 shows such a matrix (attributed to the
CreditMetrics Technical Document): a AAA issuer keeps its rating over one year with probability
$43.78\%$, while a CCC issuer is upgraded to B with probability $41.53\%$, stays CCC with $32.46\%$ and
defaults with $19.15\%$ \ts{24}{6}{54:32}.

\begin{coursenote}[On the migration table of slide 43]
In the widely quoted one-year matrices (S\&P, Moody's, CreditMetrics) roughly 90\% of AAA issuers keep
their rating; the AAA row on the slide, which sends more than half of the AAA issuers to AA, is unusual.
The OCW version marks the table's source as unknown; use its numbers as an illustration only.
\end{coursenote}

\begin{example}[Default probabilities from a migration matrix (not from the course)]\label{ch04:ex-credit}
With three states, investment grade (IG), high yield (HY) and default (D), take
\[
  P=\begin{pmatrix}0.97&0.025&0.005\\0.05&0.90&0.05\\0&0&1\end{pmatrix}.
\]
The cumulative default probabilities $(P^n)_{\cdot D}$ are $0.5\%,\ 3.46\%,\ 8.46\%$ for an IG issuer and
$5\%,\ 20.75\%,\ 33.87\%$ for an HY issuer at $n=1,5,10$ years; they are not $n$ times the one-year
probability, because issuers migrate between grades before defaulting. Since D is absorbing and reachable from
everywhere, the unique stationary distribution is $\bm e_D$: every issuer eventually defaults, and the interesting
quantities are transient ones.
\end{example}

\begin{finance}[Uses and caveats of rating chains (not discussed in class)]
Migration matrices drive credit-portfolio models (CreditMetrics), regulatory capital, and the pricing of
credit-risky bonds and credit derivatives (\cref{ch:counterparty,ch:quanto}). The Markov assumption is only an
approximation: recently downgraded issuers are more likely to be downgraded again (``rating momentum''),
and migration rates vary with the business cycle, so the matrix is not time-homogeneous.
\end{finance}

\begin{casestudy}[Up and down days of AAPL (slide 44, \file{MC_Example1.r})]
\textbf{Goal.} Check whether the direction of daily price moves has short memory \ts{24}{6}{55:39}.
\textbf{Data.} AAPL daily closing prices $P_t$ downloaded with \texttt{quantmod} (2013/14 to September 2024);
the chain is fitted on the last four and a half years \ts{24}{6}{58:51}.
\textbf{Method.} A day is Up if $P_t>P_{t-1}$ and Down if $P_t<P_{t-1}$. The two-day state on day $t$ is one
of UU, UD, DU, DD (moves on days $t-1$ and $t$), which gives a 4-state chain whose transitions are partly
forced: from UU one can only go to UU or UD, and so on (\cref{ch04:fig-mc}). The package
\texttt{markovchain} (Spedicato, 2017) estimates the transition matrix by empirical frequencies and draws the
transition graph \ts{24}{6}{59:56}. The same is repeated with three-day states (8 states) \ts{24}{6}{1:02:01}.
\textbf{Results} (as reported in class; no knitted output is available). Two up days tend to be followed by a
third ($\hat P(\mathrm{UU}\to\mathrm{UU})>0.5$), while after two down days a reversal is more likely
($\hat P(\mathrm{DD}\to\mathrm{DD})<0.5$) \ts{24}{6}{1:00:56}. In the 8-state chain, three down days are
followed by an up day with probability close to $0.6$; the largest probability of a down day, $0.510$,
follows a mixed state with two up days preceded by a down day \ts{24}{6}{1:05:15}. All states communicate
(the graph is strongly connected) \ts{24}{6}{1:04:05}.
\textbf{Lesson.} Discretising returns into a few states and fitting a Markov chain is a quick exploratory
tool, easy to rerun on other tickers or periods \ts{24}{6}{1:06:18}. \emph{Our remark:} with about 1100
trading days each of the 8 states is visited roughly 140 times, so an estimated probability near $0.5$ has a
standard error of about $\sqrt{0.25/140}\approx0.04$: $0.510$ is indistinguishable from $0.5$, and even
$0.6$ is only about two standard errors away, before accounting for having looked at 8 states. Note also that
the slide's definitions leave days with $P_t=P_{t-1}$ unclassified; one inequality should be
non-strict, as the lecturer points out \ts{24}{6}{55:39}.
\end{casestudy}

\begin{figure}[ht]
\centering
\begin{tikzpicture}[>=Stealth,every node/.style={font=\small},
    st/.style={circle,draw=defcol,thick,fill=defcol!6,minimum size=9mm,inner sep=0pt}]
  % machine chain
  \node[st] (W) at (0,0) {W};
  \node[st] (B) at (2.8,0) {B};
  \draw[->,thick] (W) to[bend left=25] node[above,font=\scriptsize]{$0.01$} (B);
  \draw[->,thick] (B) to[bend left=25] node[below,font=\scriptsize]{$0.8$} (W);
  \draw[->,thick] (W) to[loop above,looseness=6] node[above,font=\scriptsize]{$0.99$} (W);
  \draw[->,thick] (B) to[loop above,looseness=6] node[above,font=\scriptsize]{$0.2$} (B);
  \node[font=\footnotesize] at (1.4,-1.6) {(a) machine chain, \cref{ch04:ex-machine}};
  % two-day chain
  \begin{scope}[xshift=6.3cm,yshift=0.9cm]
  \node[st] (UU) at (0,0) {UU};
  \node[st] (UD) at (3,0) {UD};
  \node[st] (DU) at (0,-2.2) {DU};
  \node[st] (DD) at (3,-2.2) {DD};
  \draw[->,thick,intcol] (UU) to[loop left,looseness=6] (UU);
  \draw[->,thick,thmcol] (UU) -- (UD);
  \draw[->,thick,intcol] (UD) to[bend left=12] (DU);
  \draw[->,thick,thmcol] (UD) -- (DD);
  \draw[->,thick,intcol] (DU) -- (UU);
  \draw[->,thick,thmcol] (DU) to[bend left=12] (UD);
  \draw[->,thick,intcol] (DD) -- (DU);
  \draw[->,thick,thmcol] (DD) to[loop right,looseness=6] (DD);
  \node[font=\footnotesize] at (1.5,-3.4) {(b) two-day up/down states};
  \end{scope}
\end{tikzpicture}
\caption{(a) The two-state machine chain of the 2013 lecture. (b) Allowed transitions of the two-day
states of the AAPL case study: the second letter of today's state is the first letter of tomorrow's; green
arrows are followed by an up day, red ones by a down day. The estimated probabilities were shown in class
\ts{24}{6}{1:00:56}.}\label{ch04:fig-mc}
\end{figure}

\begin{exercise}[Time to default (not from the course)]
For the rating chain of \cref{ch04:ex-credit}, let $Q$ be the $2\times2$ block of transitions among the non-default
states. Show that the expected number of years before default, starting from IG or HY, is the vector
$(I-Q)^{-1}\bm 1$ (condition on the first step), and compute it.
\end{exercise}

% ------------------------------------------------------------------
\section{Martingales}\label{ch04:sec-mart}

\subsection{Information and conditional expectation}
Let $X_1,X_2,\dots$ be the observed process. The \emph{information set} $\mathcal F_n$ collects what is known at
time $n$, i.e.\ the values of $X_1,\dots,X_n$ (formally, the $\sigma$-algebra they generate) \ts{24}{6}{8:32}.
A random variable $Y$ is \emph{known at time $n$}, written $Y\in\mathcal F_n$, iff $Y=f(X_1,\dots,X_n)$ for
some function $f$; the increasing family $\{\mathcal F_n\}$ is a \emph{filtration}, and one writes
$\E[Z\mid\mathcal F_n]=\E[Z\mid X_1,\dots,X_n]$ \ts{24}{6}{13:00}. We use the following properties of
conditional expectation (\cref{ch:probability}), valid for integrable $Z$:
\begin{itemize}
  \item \emph{taking out what is known}: $\E[YZ\mid\mathcal F_n]=Y\,\E[Z\mid\mathcal F_n]$ if $Y\in\mathcal F_n$;
  \item \emph{tower property}: $\E\bigl[\E[Z\mid\mathcal F_n]\mid\mathcal F_m\bigr]=\E[Z\mid\mathcal F_m]$ for
        $m\le n$; in particular $\E\bigl[\E[Z\mid\mathcal F_n]\bigr]=\E[Z]$;
  \item \emph{independence}: $\E[Z\mid\mathcal F_n]=\E[Z]$ if $Z$ is independent of $X_1,\dots,X_n$.
\end{itemize}

\subsection{Definition and first properties}
\begin{definition}[Martingale]\label{ch04:def-mart}
A sequence $M_n=f_n(X_1,\dots,X_n)$, $n\ge1$, with $M_0$ a constant, is a \emph{martingale} with respect
to $\{X_n\}$ (or to $\{\mathcal F_n\}$) if for all $n\ge1$
\begin{equation}\label{ch04:eq-mart}
  \E[|M_n|]<\infty\qquad\text{and}\qquad \E[M_n\mid\mathcal F_{n-1}]=M_{n-1}
\end{equation}
\ts{24}{5}{1:13:28}. It is a \emph{submartingale} if $\E[M_n\mid\mathcal F_{n-1}]\ge M_{n-1}$ and a
\emph{supermartingale} if $\le$ (terms not used in class).
\end{definition}
Standing at time $n-1$, the best prediction of the next value is the current value: the conditional
expectations of the future are ``flat'' \ts{24}{6}{2:05}. If $M_n$ is the balance of a gambler, the game is
\emph{fair}: whatever happened so far, the expected gain of the next round is zero \ts{13}{5}{1:01:24}. The
balance of a roulette player is a supermartingale, not a martingale: the game is designed so that
$\E[M_{n+1}\mid\mathcal F_n]<M_n$ \ts{13}{5}{1:02:24}.

\begin{coursenote}[2013 versus 2024 definitions]
The 2013 notes define a martingale with respect to its own past, $X_t=\E[X_{t+1}\mid X_0,\dots,X_t]$, and
omit the integrability condition. The 2024 definition allows the conditioning information to be generated
by a different, ``driving'' sequence $X_n$ (as in the examples below). A martingale with respect to
$\{X_n\}$ is also a martingale with respect to its own past, by the tower property, since
$M_0,\dots,M_{n-1}$ are functions of $X_1,\dots,X_{n-1}$.
\end{coursenote}

\begin{proposition}\label{ch04:prop-mart}
If $\{M_n\}$ is a martingale then $\E[M_n\mid\mathcal F_m]=M_m$ for all $m\le n$, and $\E[M_n]=M_0$ for all
$n$. The increments $D_k=M_k-M_{k-1}$ satisfy $\E[D_k\mid\mathcal F_{k-1}]=0$; if they are square-integrable
they are uncorrelated, and $\Var(M_n)=\sum_{k=1}^n\E[D_k^2]$.
\end{proposition}
\begin{proof}
By the tower property, $\E[M_n\mid\mathcal F_m]=\E\bigl[\E[M_n\mid\mathcal F_{n-1}]\mid\mathcal F_m\bigr]
=\E[M_{n-1}\mid\mathcal F_m]$, and induction on $n-m$ (2013 notes, Proposition~4.2). Taking $m=0$ gives
$\E[M_n]=M_0$. For $j<k$, $D_j\in\mathcal F_{k-1}$, so
$\E[D_jD_k]=\E\bigl[D_j\,\E[D_k\mid\mathcal F_{k-1}]\bigr]=0$; hence the variance of the sum is the sum of
the variances.
\end{proof}

\subsection{Examples}\label{ch04:sec-mart-ex}
The 2024 slides give four families; each one is used later.

\begin{example}[Sums: random walks {\ts{24}{5}{1:13:28}}]\label{ch04:ex-m1}
If $X_n$ are independent with $\E[X_n]=0$, then $S_n=X_1+\dots+X_n$ is a martingale:
$\E[S_n\mid\mathcal F_{n-1}]=S_{n-1}+\E[X_n]=S_{n-1}$. In particular the SRW is a martingale
\ts{13}{5}{1:01:24}. With non-zero means, the \emph{compensated} walk $S_n-\sum_{k\le n}\E[X_k]$ is a
martingale.
\end{example}

\begin{example}[Squares: $S_n^2-n\sigma^2$ {\ts{24}{5}{1:17:44}}]\label{ch04:ex-m2}
If moreover $\Var(X_n)=\sigma^2$, then $M_n=S_n^2-n\sigma^2$ is a martingale:
\[
  \E[S_n^2\mid\mathcal F_{n-1}]=S_{n-1}^2+2S_{n-1}\E[X_n]+\E[X_n^2]=S_{n-1}^2+\sigma^2 .
\]
Its paths look very different from those of $S_n$ (slides 10--12): bounded below by $-n\sigma^2$, unbounded
above, with a long right tail \ts{24}{6}{3:09}. It measures how much the walk has spread compared with its
expected spread; its continuous-time analogue $B_t^2-t$ encodes the quadratic variation of Brownian
motion, the origin of the extra term in It\^o's formula (\cref{ch:stochcalc}).
\end{example}

\begin{example}[Products {\ts{24}{6}{4:15}}]\label{ch04:ex-m3}
If $X_n\ge0$ are independent with $\E[X_n]=1$ and $M_0=1$, then $M_n=X_1X_2\cdots X_n$ is a martingale:
$\E[M_n\mid\mathcal F_{n-1}]=M_{n-1}\E[X_n]=M_{n-1}$. The 2024 slides take $X_n\in\{0.5,\,1.5\}$ with
probability $1/2$ each: the wealth of someone who stakes everything, every period, on a fair bet that gains
or loses 50\% \ts{24}{6}{5:20}. The 2013 lecture takes $X_n=2$ with probability $1/3$ and $X_n=\tfrac12$
with probability $2/3$, again with mean $1$ \ts{13}{5}{1:03:24}.
\end{example}

The simulations in class show a few paths growing to 5--10 times the initial wealth, while most lose
almost everything \ts{24}{6}{5:20}. The calculation behind the pictures (ours): $\E[M_n]=1$ for every $n$, but
$\log M_n=\sum_k\log X_k$ is a random walk with drift $\E[\log X]=\tfrac12\ln0.75=-0.144$ per step, so
$M_n\to0$ almost surely by the law of large numbers. After 50 steps the median wealth is
$0.75^{25}\approx7.5\times10^{-4}$, and $\Prob(M_{50}\ge1)=\Prob(\mathrm{Bin}(50,\tfrac12)\ge32)\approx0.032$;
the mean $1$ is carried by rare paths (the largest possible value is $1.5^{50}\approx6\times10^8$). In the 2013
example $\E[\log X]=\tfrac13\ln2-\tfrac23\ln2<0$ as well.

\begin{intuition}[Mean versus typical behaviour]
For multiplicative processes the mean and the typical value part ways: $\E[M_n]$ is dominated by
exponentially rare paths, much as a partition function can be dominated by rare configurations. The
typical growth rate is $\E[\log X]$, not $\log\E[X]$; by Jensen's inequality $\E[\log X]<\log\E[X]=0$ unless $X$ is
constant. In finance this is the gap between arithmetic and geometric mean returns (``volatility drag''),
and the reason for the debate mentioned in class \ts{24}{5}{49:47} on whether an investor should maximise
the expected \emph{logarithm} of wealth (the Kelly criterion, of which Paul Samuelson was the best-known
critic).
\end{intuition}

\begin{example}[Exponential martingales {\ts{24}{6}{6:28}}]\label{ch04:ex-m4}
Let $Y_n$ be i.i.d.\ with moment generating function $\phi(\lambda)=\E[e^{\lambda Y}]<\infty$. Then
$X_n=e^{\lambda Y_n}/\phi(\lambda)$ are independent with mean $1$, so by \cref{ch04:ex-m3}
\begin{equation}\label{ch04:eq-expmart}
  M_n=\frac{e^{\lambda S_n}}{\phi(\lambda)^n},\qquad S_n=Y_1+\dots+Y_n,
\end{equation}
is a martingale for every fixed $\lambda$ \ts{24}{6}{7:30}. If $\lambda_0\ne0$ solves $\phi(\lambda_0)=1$,
then $M_n=e^{\lambda_0S_n}$ is a martingale. Since $\phi$ is convex with $\phi(0)=1$ and $\phi'(0)=\E[Y]$, such
a root exists (with the sign opposite to $\E[Y]$) as soon as $\E[Y]\ne0$ and $\phi$ grows beyond $1$ on that side,
e.g.\ for bounded steps taking both signs.
\end{example}
For $\pm1$ steps, $\phi(\lambda)=pe^\lambda+qe^{-\lambda}$ and $\phi(\lambda)=1$ is a quadratic equation in
$e^\lambda$ with roots $1$ and $q/p$: hence $(q/p)^{S_n}$ is a martingale (\cref{ch04:fig-ruin}, left)
\ts{24}{6}{33:19}. This is the key to the biased gambler's ruin (\cref{ch04:sec-ruin}). The exponent
$\lambda_0$ is the ``adjustment coefficient'' of actuarial ruin theory, and \eqref{ch04:eq-expmart} is the
discrete ancestor of the exponential martingale $e^{\lambda B_t-\lambda^2t/2}$ of Brownian motion and of
the change of measure behind risk-neutral pricing (\cref{ch:bs}).

\paragraph{Markov versus martingale.} The two notions are independent \ts{13}{5}{1:05:32}. The SRW is both.
A biased walk is Markov but not a martingale. A martingale that is not Markov (our example): toss a fair
coin $\xi_1$ and let $M_n=\xi_1+c(\xi_1)\sum_{k=2}^n\xi_k$ with further fair signs $\xi_k$ and step size
$c=1$ if $\xi_1=+1$, $c=2$ if $\xi_1=-1$. It is a martingale (a martingale transform, see below), but
$M_3=-1$ can be reached with either step size, so the law of $M_4$ given $M_3=-1$ depends on the earlier path.
Processes of neither kind are the rule rather than the exception.

\begin{remark}[Prices or log-prices?]\label{ch04:rem-logprice}
If $\log P_t$ is a random walk with i.i.d.\ steps $X_t$ (2024 PS3 Problem~3), the price itself satisfies
$\E[P_t\mid\mathcal F_{t-1}]=P_{t-1}\E[e^{X}]$, so $P_t$ is a martingale iff $\E[e^X]=1$. For normal steps
$\E[e^X]=e^{\mu+\sigma^2/2}$, i.e.\ iff $\mu=-\sigma^2/2$: a martingale log-price ($\mu=0$) gives an
\emph{upward}-drifting price. This lognormal correction (2024 PS2 Problem~2, \cref{ch:probability}) is the same
$-\sigma^2/2$ that appears in the drift of $\log S_t$ under Black--Scholes.
\end{remark}

% ------------------------------------------------------------------
\section{Martingale transforms and stopping times}\label{ch04:sec-stopping}

\subsection{Trading strategies as martingale transforms}
\begin{definition}[Non-anticipating sequence, martingale transform]
A sequence $\{A_n,\ n\ge1\}$ is \emph{non-anticipating} (or predictable) with respect to $\{\mathcal F_n\}$ if
$A_n\in\mathcal F_{n-1}$ for every $n$: $A_n$ is decided with the information available at time $n-1$
\ts{24}{6}{9:44}. The \emph{martingale transform} of a martingale $\{M_n\}$ by $\{A_n\}$ is
\[
  \tilde M_n=M_0+A_1(M_1-M_0)+A_2(M_2-M_1)+\dots+A_n(M_n-M_{n-1})
  =M_0+\sum_{k=1}^nA_k\,\Delta M_k
\]
\ts{24}{6}{10:51}.
\end{definition}

\begin{theorem}[Martingale transform theorem]\label{ch04:thm-transform}
If $\{M_n\}$ is a martingale and $\{A_n\}$ is non-anticipating and bounded, $|A_n|\le C_n$ for constants
$C_n$, then $\{\tilde M_n\}$ is a martingale with respect to $\{\mathcal F_n\}$ \ts{24}{6}{13:00}.
\end{theorem}
\begin{proof}
$\tilde M_n$ is a function of $X_1,\dots,X_n$ and
$\E[|\tilde M_n|]\le|M_0|+\sum_{k\le n}C_k\bigl(\E[|M_k|]+\E[|M_{k-1}|]\bigr)<\infty$. Taking out what is known,
\[
  \E[\tilde M_n-\tilde M_{n-1}\mid\mathcal F_{n-1}]=A_n\,\E[M_n-M_{n-1}\mid\mathcal F_{n-1}]=A_n\cdot0=0 .
  \qedhere
\]
\end{proof}

\begin{finance}[You cannot beat a fair game]
Read $M_n$ as the (discounted) price of an asset and $A_k$ as the number of shares held from $k-1$ to $k$,
chosen using only the information available at $k-1$ \ts{24}{6}{10:51}. Then $\tilde M_n-M_0=\sum_kA_k\Delta M_k$ is
the cumulative trading gain: exactly the ``holdings times price increment, holdings decided first'' sum of
\eqref{eq:pnl} and the discrete version of the It\^o integral $\int\theta_t\dd S_t$ (\cref{ch:stochcalc}).
The theorem says that if prices are martingales, \emph{every} bounded trading strategy has zero expected gain.
This is the content of the fundamental theorem of asset pricing read backwards: absence of arbitrage is
equivalent to the existence of a measure $\Q$ under which discounted prices are martingales; in one period
this is the relation $P_0=\beta\,\E_Q[P_T]$ of \cref{thm:ftap}.
\end{finance}

\subsection{Stopping times}
\begin{definition}[Stopping time]\label{ch04:def-stop}
A random variable $\tau$ with values in $\{0,1,2,\dots\}\cup\{\infty\}$ is a \emph{stopping time} with respect to
$\{\mathcal F_n\}$ if $\{\tau\le n\}\in\mathcal F_n$ for every $n$ (equivalently $\{\tau=n\}\in\mathcal F_n$)
\ts{24}{6}{14:04}.
\end{definition}
Viewing $\omega$ as a path, the event $\{\omega:\tau(\omega)\le n\}$ is decided by the first $n$ steps of the path:
at every time we know whether we have already stopped \ts{24}{6}{16:16}. A stopping time is a rule to quit a
game that looks only at what has happened so far \ts{13}{5}{1:09:40}.

\begin{example}[Stopping times and non-stopping times {\ts{13}{5}{1:10:49}}]\label{ch04:ex-stop}
In the coin-toss game with balance $S_n$:
\begin{enumerate}[(a)]
  \item the first time the balance reaches \$100, $\tau=\min\{n:S_n=100\}$, is a stopping time; so is the first
        time it reaches either \$100 or $-\$50$, and any first hitting time of a set;
  \item the time of the first peak (the first local maximum of the balance) is \emph{not} a stopping time:
        to know that $n$ is a peak one must see $S_{n+1}<S_n$ \ts{13}{5}{1:11:50};
  \item one step after the first peak \emph{is} a stopping time, as a student remarks \ts{13}{5}{1:12:51}.
\end{enumerate}
\end{example}

\begin{definition}[Stopped and truncated processes]
For a stopping time $\tau$, the value at the stopping time is $X_\tau=\sum_{n\ge0}X_n\,\ind\{\tau=n\}$ (on
$\{\tau<\infty\}$) \ts{24}{6}{17:24}. Since $\tau$ can be infinite, one works with the truncated times
$n\wedge\tau=\min(n,\tau)$ \ts{24}{6}{18:28} and the \emph{stopped process} $\{X_{n\wedge\tau},\ n\ge0\}$, which
follows $X$ until $\tau$ and then stays frozen. $X_{n\wedge\tau}=\sum_{k<n}X_k\ind\{\tau=k\}+X_n\ind\{\tau\ge n\}$ is a
function of $X_0,\dots,X_n$.
\end{definition}

\begin{theorem}[Stopped martingales are martingales]\label{ch04:thm-stopped}
If $\{M_n\}$ is a martingale and $\tau$ a stopping time, then $\{M_{n\wedge\tau}\}$ is a martingale; in
particular $\E[M_{n\wedge\tau}]=M_0$ for every $n$.
\end{theorem}
\begin{proof}
Let $A_k=\ind\{\tau\ge k\}=1-\ind\{\tau\le k-1\}$. It is bounded and $A_k\in\mathcal F_{k-1}$, because
$\{\tau\le k-1\}\in\mathcal F_{k-1}$. The transform telescopes up to the stopping time:
\[
  M_0+\sum_{k=1}^nA_k(M_k-M_{k-1})=M_0+\sum_{k=1}^{n\wedge\tau}(M_k-M_{k-1})=M_{n\wedge\tau},
\]
which is a martingale by \cref{ch04:thm-transform} \ts{24}{6}{20:29}. Then apply \cref{ch04:prop-mart}.
\end{proof}

\begin{coursenote}[Erratum: slides 17, 20 and 23 of \file{lec05}]
\begin{itemize}
  \item Slide 23 (proof above): ``else replace $M_n$ by $M_n'=M_n-M_{n-1}$'' should read $M_n'=M_n-M_0$ (to reduce to
        $M_0=0$ one subtracts the initial value, not the previous one). In the displayed identity the right-hand side
        should have $n$ in place of the summation index $k$:
        $\sum_{k=1}^nA_k(M_k-M_{k-1})=M_\tau\ind(\tau\le n-1)+M_n\ind(\tau\ge n)=M_{n\wedge\tau}$.
  \item Slide 20 calls $\{0,1,2,\dots\}$ the ``state space'' of the process; in the definition of a stopping time it
        plays the role of the \emph{time index set}, in which $\tau$ takes values (together with $\infty$).
  \item Slide 17: in the special case the sum should read $S_n=\sum_{i=1}^nY_i$, and $\lambda_0$ must be a root of
        $\phi(\lambda_0)=1$ other than the trivial root $\lambda_0=0$.
\end{itemize}
\end{coursenote}

\begin{exercise}[Stopping times (not from the course)]
Let $\sigma,\tau$ be stopping times. Show that $\sigma\wedge\tau$, $\sigma\vee\tau$ and $\tau+1$ are stopping times, but that
$\tau-1$ in general is not. Deduce that ``one step after the first peak'' is a stopping time and ``the first peak'' is
not (\cref{ch04:ex-stop}).
\end{exercise}

\subsection{The optional stopping theorem}
\cref{ch04:thm-stopped} gives $\E[M_{n\wedge\tau}]=M_0$ at every fixed $n$. The \emph{optional stopping theorem}
passes to the limit $n\to\infty$, where $M_{n\wedge\tau}\to M_\tau$ on $\{\tau<\infty\}$.

\begin{theorem}[Optional stopping]\label{ch04:thm-ost}
Let $\{M_n\}$ be a martingale and $\tau$ a stopping time. Then $\E[M_\tau]=\E[M_0]$ in each of the following
cases:
\begin{enumerate}[(i)]
  \item $\tau\le T$ for a constant $T$ (2013 notes, Theorem~5.3, ``Doob's optional stopping theorem, weak form'')
        \ts{13}{5}{1:14:06};
  \item $\tau<\infty$ almost surely and the stopped process is bounded, $|M_{n\wedge\tau}|\le K$ for all $n$;
  \item $\E[\tau]<\infty$ and the increments are bounded, $|M_n-M_{n-1}|\le c$.
\end{enumerate}
\end{theorem}
\begin{proof}
(i) For $n=T$, $M_{T\wedge\tau}=M_\tau$, so $\E[M_\tau]=\E[M_0]$ by \cref{ch04:thm-stopped}. (This is the 2013 proof:
$M_\tau=M_0+\sum_{i=0}^{T-1}(M_{i+1}-M_i)\ind\{\tau\ge i+1\}$, and each term has zero mean after conditioning on
$\mathcal F_i$, because $\{\tau\ge i+1\}$ is decided by time $i$.)

(ii) Since $\tau<\infty$, $M_{n\wedge\tau}=M_\tau$ for all $n\ge\tau$, so $M_{n\wedge\tau}\to M_\tau$ almost surely; by
bounded convergence $\E[M_\tau]=\lim_n\E[M_{n\wedge\tau}]=\E[M_0]$.

(iii) $|M_{n\wedge\tau}|\le|M_0|+c\,\tau$, an integrable bound, so dominated convergence applies as in (ii).
\end{proof}

``A mortal being has no winning strategy'' in a fair game (2013 notes): whatever rule you use to decide when to
quit, as long as it does not peek into the future and you cannot play for an unbounded time with unbounded credit,
your expected final balance equals your initial one \ts{13}{5}{1:06:34}, \ts{13}{5}{1:15:08}. The hypotheses are
not decorations: 2013 notes, Exercise~5.4 (\cref{ch04:ex-note54}), asks why stopping the SRW when it first
reaches \$100 does not contradict the theorem although $\E[S_\tau]=100$.

\begin{finance}[The ``martingale'' betting system (not discussed in class)]
Bet \$1 on a fair coin; after each loss double the stake; stop at the first win. With $A_k=2^{k-1}$ on
$\{\tau\ge k\}$ and $\tau$ the first win, the net gain at $\tau$ is $2^{\tau-1}-(1+2+\dots+2^{\tau-2})=+1$ with
certainty, although $\E[\tau]=2<\infty$. Condition (iii) fails because the stakes, i.e.\ the increments of the
transformed process, are unbounded, and so does (ii), since the loss before the first win is unbounded. At any fixed
horizon $n$ the game is still fair: with probability $1-2^{-n}$ you have won \$1, with probability $2^{-n}$ you have lost
$2^n-1$, and the mean is $0$. This doubling scheme is the historical meaning of ``martingale'' in gambling (the
lecturer recalls the other meaning, a horse's bridle \ts{24}{5}{1:18:49}). Its trading version, doubling a losing
position to win it back, needs unlimited capital and is a classic route to ruin.
\end{finance}

\begin{exercise}[PS3 (2024), Problem 3: random walk for the log price]\label{ch04:ex-ps3-3}
Let $P_t$ be an asset price, $S_t=\log P_t$ and $X_t=S_t-S_{t-1}=\log(P_t/P_{t-1})$, with the $X_t$ i.i.d., $\E[X_t]=\mu$
and $\Var X_t=\sigma^2>0$.
\begin{enumerate}[(a)]
  \item For which $(\mu,\sigma^2)$ is $\{S_t\}$ a Markov process? Explain.
  \item For which $(\mu,\sigma^2)$ is $\{S_t\}$ a martingale? Explain.
  \item Let $S_0=0$ and $\tau=\min\{t:\ S_t>0.20\ \text{or}\ S_t<-0.20\}$, the first time the log price rises above
      $0.20$ or falls below $-0.20$. If the steps are i.i.d.\ $N(\mu,\sigma^2)$, which conditions on $\mu$ and $\sigma^2$ make
      $\Prob(S_\tau>0.20)>\tfrac12$? (The problem statement says ``$S_t$ are i.i.d.\ Normal''; the steps $X_t$ are meant.)
\end{enumerate}
\end{exercise}

\begin{exercise}[A martingale that is not Markov (not from the course)]
For the process $M_n=\xi_1+c(\xi_1)\sum_{k=2}^n\xi_k$ of \cref{ch04:sec-mart-ex}, write $M_n$ as a martingale
transform, and compute $\Prob(M_4=-3\mid M_3=-1)$ and $\Prob(M_4=-3\mid M_3=-1,\ M_1=-1)$ to show that $\{M_n\}$ is not
Markov.
\end{exercise}

\begin{exercise}[\file{lecnote5}, Exercise 5.4]\label{ch04:ex-note54}
In the coin-toss game the gambler stops the first time her balance reaches \$100. By construction $\E[S_\tau]=100$. Does
this contradict the optional stopping theorem? Which hypothesis of \cref{ch04:thm-ost} fails?
\end{exercise}

% ------------------------------------------------------------------
\section{Gambler's ruin}\label{ch04:sec-ruin}
Two players bet \$1 per round on a fair coin until one of them is broke. Let $S_n$ be my net gain after $n$
rounds, with my bankroll $B$ and my opponent's bankroll $A$ ($A,B$ positive integers) \ts{24}{6}{21:35},
\ts{13}{5}{23:03}. The game ends at
\[
  \tau=\min\{n:\ S_n=+A\ \text{or}\ S_n=-B\},
\]
a stopping time (a first hitting time), and I take all my opponent's money iff $S_\tau=+A$ \ts{24}{6}{22:46}.

\begin{lemma}\label{ch04:lem-finite}
For a walk with steps $\pm1$ and $\Prob(X=+1)=p\in(0,1)$, $\Prob(\tau>n(A+B))\le(1-p^{A+B})^n$. Hence
$\tau<\infty$ almost surely and $\E[\tau]<\infty$.
\end{lemma}
\begin{proof}
Split time into blocks of $A+B$ steps. If all steps of a block are $+1$ the walk, wherever it is in $(-B,A)$,
leaves the interval during that block. Blocks are independent and each has only $+1$ steps with probability
$p^{A+B}$, so surviving $n$ blocks has probability at most $(1-p^{A+B})^n$. Then
$\E[\tau]=\sum_{m\ge0}\Prob(\tau>m)\le(A+B)\sum_{n\ge0}(1-p^{A+B})^n<\infty$.
\end{proof}

\begin{proposition}[Fair gambler's ruin]\label{ch04:prop-ruin}
For the SRW,
\[
  \Prob(S_\tau=+A)=\frac{B}{A+B},\qquad \Prob(S_\tau=-B)=\frac{A}{A+B},\qquad \E[\tau]=AB .
\]
\end{proposition}
\begin{proof}
$S_n$ is a martingale and $|S_{n\wedge\tau}|\le\max(A,B)$, so by \cref{ch04:lem-finite} and optional stopping (b),
$0=\E[S_\tau]=A\,\Prob(S_\tau=A)-B\,\bigl(1-\Prob(S_\tau=A)\bigr)$ \ts{24}{6}{23:52}, \ts{24}{6}{24:54}, which gives the
first formula. For the second, $M_n=S_n^2-n$ is a martingale (\cref{ch04:ex-m2} with $\sigma^2=1$), so by
\cref{ch04:thm-stopped}
\[
  \E\bigl[S_{n\wedge\tau}^2\bigr]=\E[n\wedge\tau]\qquad\text{for every }n.
\]
As $n\to\infty$ the left side tends to $\E[S_\tau^2]$ by bounded convergence ($S^2_{n\wedge\tau}\le\max(A^2,B^2)$) and the
right side to $\E[\tau]$ by monotone convergence ($n\wedge\tau\uparrow\tau$). Hence \ts{24}{6}{29:13}
\[
  \E[\tau]=\E[S_\tau^2]=A^2\,\frac{B}{A+B}+B^2\,\frac{A}{A+B}=AB .\qedhere
\]
\end{proof}

With equal bankrolls each player wins with probability $\tfrac12$, by symmetry \ts{13}{5}{25:08}. Playing until I
win \$100 or lose \$50 ($A=100$, $B=50$), I win with probability $1/3$ \ts{13}{5}{26:08}, and the game lasts
$5000$ rounds on average. In a fair game the richer player is more likely to win everything
\ts{24}{6}{27:05}, and the expected duration is the product of the bankrolls \ts{24}{6}{29:13}.

\begin{coursenote}[Erratum: slides 24--26 of \file{lec05}]
\begin{itemize}
  \item ``Problem: Solve for $\Prob(\tau=+A)$'' (and the title of slide 25) should read $\Prob(S_\tau=+A)$: $\tau$ is a
        time, $S_\tau$ the level hit. A student pointed this out in class and the lecturer confirmed \ts{24}{6}{26:05}.
  \item ``$\lim_{n\to\infty}\E[S_{n\wedge\tau}]=S_\tau$ w.p.~1'' mixes a number and a random variable: the correct
        statement is $S_{n\wedge\tau}\to S_\tau$ almost surely (because $\Prob(\tau=\infty)=0$), and then
        $\E[S_\tau]=\lim_n\E[S_{n\wedge\tau}]=0$ by bounded convergence.
  \item Slide 26: ``$M_{n\wedge\tau}\le\max(A^2,B^2)+\tau$, bounded r.v.'' is not a valid justification, since $\tau$ is not
        bounded. The limit must be taken separately in the two terms of $M_{n\wedge\tau}=S^2_{n\wedge\tau}-n\wedge\tau$,
        with bounded and monotone convergence, as in the proof above.
\end{itemize}
\end{coursenote}

\paragraph{The 2013 route: first-step analysis.} Let $f(k)$ be the probability of hitting $+A$ before $-B$ when
starting from $k\in\{-B,\dots,A\}$ \ts{13}{5}{29:19}. Conditioning on the first toss and using the stationarity of
the increments,
\[
  f(k)=\tfrac12f(k+1)+\tfrac12f(k-1)\quad(-B<k<A),\qquad f(A)=1,\quad f(-B)=0
\]
\ts{13}{5}{30:20}. The increments $f(k+1)-f(k)$ are therefore constant: $f(-B+r)=\alpha r$ with
$\alpha=1/(A+B)$, and $f(0)=B/(A+B)$ (2013 notes, Section~2).

\begin{intuition}[Harmonic functions and electrostatics]
The recursion says that $f$ is \emph{discrete-harmonic}: its value is the average of its neighbours, i.e.\ its discrete
Laplacian vanishes. Equivalently $f(S_{n\wedge\tau})$ is a martingale, which is why the two methods agree. The hitting
probability is the solution of a Dirichlet problem, like the potential between two plates held at potentials $0$ and
$1$, hence linear. In continuous time the same idea links expectations of Brownian functionals to PDEs (backward
Kolmogorov and Feynman--Kac equations, used to price options in \cref{ch:bs}).
\end{intuition}

\begin{remark}[A student's symmetry argument {\ts{13}{5}{27:11}}]
For $A=100$, $B=50$ a student guessed $1/3$ as follows. Let $g(a,b)$ be the probability of gaining $a$ before losing
$b$. To gain $100$ before losing $50$, the walk must first reach $+50$ before $-50$ (probability $\tfrac12$); from $+50$
it must then gain $50$ before losing $100$, with probability $g(50,100)=1-g(100,50)$ by the reflection
$S\mapsto-S$. So $g=\tfrac12(1-g)$ and $g=1/3$. The lecturer was unsure it could be made rigorous; it can: restarting the
walk at a stopping time gives again a SRW independent of the past (the \emph{strong Markov property}, not proved in
this course).
\end{remark}

\paragraph{The biased walk.} Now $\Prob(X=+1)=p$, $\Prob(X=-1)=q=1-p$, $p\ne\tfrac12$ \ts{24}{6}{30:16}. First-step
analysis still works but is laborious \ts{24}{6}{31:16}; the exponential martingale solves the problem in two lines
\ts{24}{6}{32:19}.

\begin{proposition}[Biased gambler's ruin]\label{ch04:prop-biased}
With $r=q/p\ne1$,
\[
  \Prob(S_\tau=+A)=\frac{r^B-1}{r^{A+B}-1}.
\]
\end{proposition}
\begin{proof}
By \cref{ch04:ex-m4}, $M_n=r^{S_n}$ is a martingale with $M_0=1$. Since $-B\le S_{n\wedge\tau}\le A$, $M_{n\wedge\tau}$ lies
between $r^{-B}$ and $r^{A}$, and optional stopping (b) with \cref{ch04:lem-finite} gives \ts{24}{6}{34:22}
\[
  1=\E[r^{S_\tau}]=r^{A}\,\Prob(S_\tau=A)+r^{-B}\bigl(1-\Prob(S_\tau=A)\bigr)
  \quad\Longrightarrow\quad \Prob(S_\tau=A)=\frac{1-r^{-B}}{r^A-r^{-B}}=\frac{r^B-1}{r^{A+B}-1}.\qedhere
\]
\end{proof}
As $p\to\tfrac12$, $r=1+\varepsilon$ with $\varepsilon\to0$, and the ratio tends to $B\varepsilon/((A+B)\varepsilon)=B/(A+B)$,
recovering the fair case. The expected duration (not derived in class) follows from the martingale
$S_n-n(p-q)$ (\cref{ch04:ex-m1}) and optional stopping (c) (bounded increments, $\E[\tau]<\infty$):
$\E[S_\tau]=(p-q)\,\E[\tau]$ (Wald's identity), so
\[
  \E[\tau]=\frac{A\,\Prob(S_\tau=A)-B\,\Prob(S_\tau=-B)}{p-q}.
\]

\begin{figure}[ht]
\centering
\begin{tikzpicture}
\begin{axis}[name=phiax,width=0.47\linewidth,height=5.2cm,domain=-2:1,samples=100,ymin=0,ymax=3,xmin=-2,xmax=1,
  xlabel={$\lambda$},ylabel={$\phi(\lambda)$},title={\small MGF of a step, $p=0.8$},
  tick label style={font=\scriptsize},label style={font=\small}]
  \addplot[defcol,thick] {0.8*exp(x)+0.2*exp(-x)};
  \addplot[black!40,dashed] {1};
  \draw[thmcol,thick] (axis cs:-1.3863,0) -- (axis cs:-1.3863,3);
  \node[thmcol,font=\scriptsize,anchor=south west,align=left] at (axis cs:-1.35,1.6) {$\lambda_0=\ln(q/p)$\\$=-1.386$};
\end{axis}
\begin{axis}[at={(phiax.east)},anchor=west,xshift=1.4cm,width=0.47\linewidth,height=5.2cm,
  domain=0:40,samples=41,ymin=0,ymax=1,xmin=0,xmax=40,
  xlabel={my bankroll $B$ (with $A+B=40$)},ylabel={$\Prob(S_\tau=+A)$},title={\small Probability of winning all},
  tick label style={font=\scriptsize},label style={font=\small},
  legend style={font=\scriptsize,at={(0.03,0.97)},anchor=north west,draw=none,fill=none}]
  \addplot[intcol,thick] {((0.47/0.53)^x-1)/((0.47/0.53)^40-1)};
  \addplot[intcol!60,thick,dashed] {((0.49/0.51)^x-1)/((0.49/0.51)^40-1)};
  \addplot[black,thick] {x/40};
  \addplot[thmcol!60,thick,dashed] {((0.51/0.49)^x-1)/((0.51/0.49)^40-1)};
  \addplot[thmcol,thick] {((0.53/0.47)^x-1)/((0.53/0.47)^40-1)};
  \legend{$p=0.53$,$p=0.51$,$p=0.50$,$p=0.49$,$p=0.47$}
\end{axis}
\end{tikzpicture}
\caption{Left: the moment generating function $\phi(\lambda)=pe^\lambda+qe^{-\lambda}$ of a $\pm1$ step equals $1$ at
$\lambda=0$ and at $\lambda_0=\ln(q/p)$, as on slide 28 \ts{24}{6}{33:19}. Right: probability of winning the
opponent's whole bankroll, \cref{ch04:prop-biased}, as a function of my initial capital when the total capital is
$40$. A bias of a few percent per bet turns the straight line of the fair game into an S-shaped curve.}\label{ch04:fig-ruin}
\end{figure}

\begin{casestudy}[Probability surface in R (slides 30--31)]
The slides define \texttt{fcn.probA(A,B)} implementing \cref{ch04:prop-biased} with $p=0.53$, evaluate it on the grid
$A,B\in\{1,\dots,30\}$ with \texttt{outer} (after \texttt{Vectorize}), and draw the surface (\texttt{persp3D}) and its
contour lines \ts{24}{6}{35:30}. Some values: $A=B=10$ gives $0.769$ and $A=B=20$ gives $0.917$. Against an opponent of
unlimited wealth ($A\to\infty$) the probability of never being ruined is $1-r^B$: $0.45$ for $B=5$, $0.70$ for $B=10$ and
$0.97$ for $B=30$. This is why the contour lines on slide~31 become horizontal for large $A$.
\end{casestudy}

\begin{finance}[Edges and ruin (numbers ours)]
The limits $A\to\infty$ make the 2013 remark precise: ``if one has some advantage over an opponent in some game, then no
matter how small that advantage is, he/she will win in the long run''. With $p>q$ the probability of never being ruined
against an infinitely rich opponent is $1-(q/p)^B>0$; with $p<q$ ruin is certain. The casino is the infinitely rich
opponent with the edge: betting \$1 on red at European roulette ($p=18/37$) with \$100, aiming to reach \$200, succeeds
with probability $0.0045$ after about $3700$ bets on average; aiming for \$20 starting from \$10 succeeds with probability
$0.37$. For a trader, $p-q$ is the edge per trade and $B$ the capital in units of the bet size: a small edge protects
against ruin only if the bets are small relative to the capital.
\end{finance}

\begin{coursenote}[2013 versus 2024 treatment]
2013 derives the fair ruin probability by first-step analysis and then, at the end of the lecture, re-derives
$\Prob=1/3$ for $(100,-50)$ as an application of optional stopping, glossing over the fact that the hitting time is
unbounded \ts{13}{5}{1:16:09}. 2024 builds the machinery first (transforms, stopping times, stopped martingales), then
computes the ruin probability, the expected duration and the biased case, all by optional stopping. 2013 has instead
the link between stationary distributions and Perron--Frobenius, which 2024 does not discuss.
\end{coursenote}

\begin{exercise}[\file{lecnote5}, Exercise 5.5]
For positive integers $a,b$ the player stops the first time the balance equals $a$ or $-b$. Find the distribution of
$S_\tau$. (This is \cref{ch04:prop-ruin}; redo it without looking. The notes write ``$X_\tau=b$'' for the second
outcome; $-b$ is meant.)
\end{exercise}

\begin{exercise}[Biased ruin by first-step analysis (not from the course)]
For the biased walk let $f(k)=\Prob(\text{hit }A\text{ before }-B\mid S_0=k)$. Show that
$f(k)=pf(k+1)+qf(k-1)$, solve this linear recursion (characteristic roots $1$ and $q/p$) with $f(A)=1$, $f(-B)=0$, and
recover \cref{ch04:prop-biased}. Then use Wald's identity to show that the expected duration tends to $AB$ as
$p\to\frac12$.
\end{exercise}

% ------------------------------------------------------------------
\section{Looking ahead}\label{ch04:sec-ahead}
\begin{itemize}
  \item \emph{Random walk $\to$ Brownian motion} (\cref{ch:sp2}): independent stationary increments and the $\sqrt t$
        scale survive the limit; $S_n^2-n$ becomes $B_t^2-t$ and $e^{\lambda S_n}/\phi(\lambda)^n$ becomes
        $e^{\lambda B_t-\lambda^2t/2}$.
  \item \emph{Martingale transforms $\to$ It\^o integrals} (\cref{ch:stochcalc,ch:sde}): gains from trading, and the
        reason why hedged portfolios under the risk-neutral measure are martingales (\cref{ch:bs}).
  \item \emph{Markov chains $\to$ Markov diffusions}: Chapman--Kolmogorov becomes the Fokker--Planck equation, and
        first-step analysis becomes the backward equation behind option-pricing PDEs.
  \item \emph{Optional stopping} returns in first-passage problems: barrier options, default times in structural credit
        models, American-style early exercise.
  \item The \emph{Poisson process}, the continuous-time counting process of events with exponential waiting times, is
        not covered in these lectures; 2013 PS2 builds the Poisson distribution from exponential waiting times
        (\cref{ch:probability}).
\end{itemize}

% ------------------------------------------------------------------
\begin{keyformulas}
\begin{align*}
&\text{Random walk:} && \E[S_n]=n\mu,\quad \Var S_n=n\sigma^2,\\
& && (S_n-n\mu)/(\sigma\sqrt n)\to N(0,1)\\
&\text{Markov chain:} && \Prob(X_0=i_0,\dots,X_n=i_n)
  =p_{i_0}P_{i_0i_1}\cdots P_{i_{n-1}i_n},\\
& && P^{(n)}=P^n,\quad \bm p^{(n)}=\bm pP^n\\
&\text{Stationary dist.:} && \bm\pi P=\bm\pi;\quad P^k>0\ \Rightarrow\
  P^{(n)}_{ij}\to\pi_j\ \text{geometrically}\\
&\text{Martingale:} && \E[M_n\mid\mathcal F_{n-1}]=M_{n-1},\\
& && \E[M_n\mid\mathcal F_m]=M_m,\quad \E[M_n]=M_0\\
&\text{Examples:} && S_n,\quad S_n^2-n\sigma^2,\\
& && \textstyle\prod_kX_k\ (\E[X]=1),\quad e^{\lambda S_n}/\phi(\lambda)^n\\
&\text{Transform:} && \tilde M_n=M_0+\textstyle\sum_{k\le n}A_k\,\Delta M_k,\\
& && A_k\in\mathcal F_{k-1}\ \Rightarrow\ \text{martingale}\\
&\text{Opt. stopping:} && \E[M_\tau]=\E[M_0]\ \text{if $\tau$ bounded,}\\
& && \text{or $M_{n\wedge\tau}$ bounded, or $\E[\tau]<\infty$ \& bounded steps}\\
&\text{Fair ruin:} && \Prob(S_\tau=A)=\tfrac{B}{A+B},\quad \E[\tau]=AB\\
&\text{Biased ruin:} && \Prob(S_\tau=A)=\tfrac{(q/p)^B-1}{(q/p)^{A+B}-1},\\
& && \E[\tau]=\tfrac{\E[S_\tau]}{p-q}
\end{align*}
\end{keyformulas}
